% Mirrors `AlgoLib.DataStructures.SegmentTree.*`
%
\subsection{Segment trees}
% Mirrors `AlgoLib.DataStructures.SegmentTree.Basic` and `AlgoLib.DataStructures.SegmentTree.Cost`

Throughout, $\alpha$ is an additive monoid; commutativity is not assumed, so sums are taken
over lists in index order. Every operation is written once, in the cost monad
\texttt{TimeM} of CSLib, charging one tick per node visited. The function one reasons
about is the value component of that computation, and the complexity bounds are statements
about its time component, so the bounds concern exactly the verified algorithm. The bounds
assume that addition in $\alpha$ takes constant time.

\begin{definition}[Segment tree]
  \label{def:segment-tree}
  \lean{SegmentTree, SegmentTree.valAt, SegmentTree.Valid}
  \leanok
  A \emph{segment tree} is a binary tree in which every subtree covers a half-open
  interval $[\ell, h)$ of indices and caches a value. A leaf covers $[\ell, \ell + 1)$. The
  tree is \emph{valid} if the children of every internal node cover consecutive intervals
  partitioning the interval of the node, and the value cached at the node is the sum of the
  values cached at its children. The tree \emph{represents} the array $a$ where $a_i$ is the
  value of the leaf covering $i$, and $a_i = 0$ outside the covered interval.
\end{definition}

\begin{definition}[Build, query, update]
  \label{def:segment-tree-ops}
  \lean{SegmentTree.build, SegmentTree.query, SegmentTree.update}
  \leanok
  \uses{def:segment-tree}
  \texttt{build} constructs a tree from a nonempty array by splitting $[\ell, h)$ at
  $\lfloor (\ell + h) / 2 \rfloor$. \texttt{query} on $[q_\ell, q_r)$ returns $0$ at a node
  disjoint from the range, the cached value at a node covered by it, and the sum of the
  answers of both children otherwise. \texttt{update} replaces the value at one index and
  recomputes the cached values on the path from that leaf to the root.
\end{definition}

\begin{theorem}[Correctness of build]
  \label{thm:segment-tree-build}
  \lean{SegmentTree.valid_build, SegmentTree.valAt_build}
  \leanok
  \uses{def:segment-tree-ops}
  The tree built from an array $a$ is valid and represents $a$.
\end{theorem}
\begin{proof}
  \leanok
  By strong induction on $h - \ell$, following the recursion of \texttt{build}.
\end{proof}

\begin{theorem}[Correctness of query]
  \label{thm:segment-tree-query}
  \lean{SegmentTree.query_eq_sum_valAt, SegmentTree.query_build}
  \leanok
  \uses{def:segment-tree-ops}
  On a valid tree covering $[\ell, h)$ and representing $a$, the query on $[q_\ell, q_r)$
  returns $\sum_{i = \max(\ell, q_\ell)}^{\min(h, q_r) - 1} a_i$.
\end{theorem}
\begin{proof}
  \leanok
  By induction on validity. At a covered node the cached value is the sum of the represented
  array over the interval of the node; at a partially covered node the range splits at the
  boundary between the two children.
\end{proof}

\begin{theorem}[Correctness of update]
  \label{thm:segment-tree-update}
  \lean{SegmentTree.Valid.update, SegmentTree.valAt_update, SegmentTree.query_update_build}
  \leanok
  \uses{def:segment-tree-ops}
  \texttt{update} preserves validity, and on a valid tree the represented array changes
  exactly at the updated index.
\end{theorem}
\begin{proof}
  \leanok
  By induction on validity, descending into the child whose interval contains the index.
\end{proof}

\begin{theorem}[Height of a built tree]
  \label{thm:segment-tree-height}
  \lean{SegmentTree.height_build, SegmentTree.height_update}
  \leanok
  \uses{def:segment-tree-ops}
  The tree built from an array of size $n$ has height exactly $\lceil \log_2 n \rceil$, and
  \texttt{update} does not change the height.
\end{theorem}
\begin{proof}
  \leanok
  The recurrence is $T(n) = \max(T(\lfloor n/2 \rfloor), T(\lceil n/2 \rceil)) + 1$, and
  the larger half dominates. The floor logarithm is not a bound: for $n = 3$ the height is
  $2$.
\end{proof}

\begin{theorem}[Cost of the operations]
  \label{thm:segment-tree-cost}
  \lean{SegmentTree.buildT_time, SegmentTree.updateT_time_le, SegmentTree.queryT_time_le, SegmentTree.updateT_build_time_le, SegmentTree.queryT_build_time_le}
  \leanok
  \uses{thm:segment-tree-height}
  On an array of size $n$, \texttt{build} visits exactly $2n - 1$ nodes. On a tree of
  height $d$, \texttt{update} visits at most $d + 1$ nodes and, if the tree is valid,
  \texttt{query} visits at most $4d + 1$ nodes. Hence both are $O(\log n)$ on a built tree.
\end{theorem}
\begin{proof}
  \leanok
  \texttt{update} follows one path. For \texttt{query}, the naive recurrence
  $T(d) \le 2T(d - 1) + 1$ is exponential. Instead, once the query range has split at a
  node, every further descent is one-sided: the range contains one end of the subtree, so
  one child is disjoint or covered and costs a single visit. One-sided descents cost at most
  $2d + 1$ visits, and combining the two sides gives $4d + 1$.
\end{proof}

\subsection{Lazy segment trees}
% Mirrors `AlgoLib.DataStructures.SegmentTree.Lazy.Basic` (work in progress)

A lazy segment tree additionally stores, at every internal node, a \emph{tag} $t \in M$
pending to be applied to its children, where the monoid $M$ acts distributively on
$\alpha$. A range update applies its tag at the covered nodes in constant time, and a tag
is pushed to the children only when an operation descends past its node. The cost bounds
for lazy segment trees are not yet formalized.

\begin{definition}[Lazy segment tree]
  \label{def:lazy-segment-tree}
  \lean{LazySegmentTree, LazySegmentTree.applyTag, LazySegmentTree.valAt, LazySegmentTree.Valid}
  \leanok
  \uses{def:segment-tree}
  The value cached at a node already includes the effect of its own tag, so validity
  requires the value at a node with tag $t$ to be $t \cdot (v_L + v_R)$, where $v_L, v_R$
  are the values cached at its children. The represented array applies every tag met on
  the way from the root to a leaf.
\end{definition}

\begin{definition}[Range traversal]
  \label{def:lazy-range-traversal}
  \lean{LazySegmentTree.rangeRecT, LazySegmentTree.query, LazySegmentTree.rangeUpdate}
  \leanok
  \uses{def:lazy-segment-tree}
  All range operations share one traversal: a node disjoint from the range or covered by it
  is answered at once, and a partially covered node pushes its tag and recurses on both
  children. \texttt{query} and \texttt{rangeUpdate} are instances of this traversal.
\end{definition}

\begin{theorem}[Correctness of lazy query and range update]
  \label{thm:lazy-segment-tree-correct}
  \lean{LazySegmentTree.query_eq_sum_valAt, LazySegmentTree.Valid.rangeUpdate, LazySegmentTree.valAt_rangeUpdate, LazySegmentTree.query_rangeUpdate_build}
  \leanok
  \uses{def:lazy-range-traversal}
  On a valid tree, \texttt{query} returns the sum of the represented array over the range.
  \texttt{rangeUpdate} on $[q_\ell, q_r)$ with tag $m$ preserves validity and acts by $m$
  on exactly the represented values with index in $[q_\ell, q_r)$.
\end{theorem}
\begin{proof}
  \leanok
  By an induction principle in which the hypotheses at a node concern its children with
  the tag of the node pushed, matching the recursive calls of the traversal.
\end{proof}
